%% P33 --- Anatomia treningu
%%
%% Bliźniak: programs/en/P33-training-run.tex. Ramka w ramkę, makro w makro;
%% tools/parity.py porównuje oba pliki.
%%
%% UWAGA DLA TŁUMACZA: C4, C8 i C12 porównują znaczniki strukturalne, wzory i
%% literały liczbowe PO KOLEI. Cyfra zostaje cyfrą, a słowo zostaje słowem;
%% odwołanie do programu prowadzi zdanie zamiast stać za wzorem. Słownictwo
%% czterech kształtów jest przejęte z polskiego F03, który przekazuje je tu z
%% nazwy: \enquote{plateau}, \enquote{skok}, \enquote{podejrzanie gładki
%% spadek}. \enquote{batch} to \enquote{partia}, jak w P21.

\program{Anatomia treningu}
\label{prog:P33}
\index{trening}

\begin{outcomes}
\outcome{1--8}{wymienić sześć elementów jednego kroku treningu i powiedzieć,
  który program tej książki dał ci każdy z nich}
\outcome{9--14}{powiedzieć, które z tych sześciu zgłaszają cokolwiek, gdy są
  zrobione źle, a które nie}
\outcome{15--21}{powiedzieć, jak długo musi trwać płaski odcinek, zanim stanie
  się dowodem, że trening przestał się uczyć}
\outcome{22--27}{powiedzieć, dlaczego długi trening zawiera duże odchylenia
  pojedynczych kroków niezależnie od tego, co jest z nim nie tak, i jaki
  rozmiar jest zaskakujący}
\outcome{28--33}{powiedzieć, co wygładzona krzywa mówi o wygładzaniu, a nie
  tylko o treningu}
\outcome{34--37}{powiedzieć, który z czterech kształtów nie potrzebuje
  żadnej statystyki}
\outcome{38--50}{dopasować prawo potęgowe, powiedzieć, co jego reszta wiąże, a
  czego nie, i powiedzieć, czego dotyczy zgłaszana przez nie niepewność}
\outcome{51--54}{powiedzieć, które twierdzenie tego programu nie zostało
  sprawdzone na prawdziwym treningu}
\end{outcomes}

\begin{quiz}
\item Krok treningu liczy stratę, różniczkuje ją i zmienia wagi. Wymień sześć
  programów tej książki, które są właścicielami tych elementów.
  \teachesat{3--4}
  \answerto{Strata należy do Programu~\ref{prog:P26}, jej gradient do
  Programu~\ref{prog:P18}, aktualizacja do Programu~\ref{prog:P20},
  arytmetyka do Programu~\ref{prog:P01}, granica rozmiaru kroku do
  Programu~\ref{prog:P17}, a próbka do Programu~\ref{prog:P21}.}
\item Ile z tych sześciu elementów zgłasza błąd, ostrzeżenie albo wyjątek,
  gdy jest zrobione źle?
  \teachesat{9--10}
  \answerto{$\val{p33.loud}$ z $\val{p33.pieces}$. Pozostałe cztery to
  poprawne obliczenie niewłaściwej wielkości, więc nie ma w nich żadnej
  niedozwolonej operacji, którą cokolwiek mogłoby wychwycić.}
\item Obserwujesz od $\val{p33.watch}$ kroków, a strata teraz nie jest niższa
  niż wtedy. Czy trening utknął?
  \teachesat{15--16}
  \answerto{Nie. Zdrowy trening wygląda dokładnie tak w około
  $\val{p33.p.flat}$ procent przypadków, a płaski odcinek staje się dowodem
  dopiero po mniej więcej $\val{p33.plateau.k}$ krokach.}
\item Trening liczący $\val{p33.steps}$ kroków zawiera jeden krok, którego
  strata leży cztery rozrzuty powyżej trendu. Ilu takich kroków się
  spodziewałeś?
  \teachesat{22--23}
  \answerto{Około $\val{p33.spike.expected}$ z nich. Długi trening zawiera
  duże odchylenia pojedynczych kroków z samej swojej konstrukcji, więc jedno
  z nich nie jest dowodem niczego.}
\item Panel wygładza krzywą straty przy $\beta = \val{p33.beta.dash}$. Jaką
  część prawdziwego rozrzutu krok po kroku pokazuje ten obraz?
  \teachesat{30--31}
  \answerto{$\val{p33.smooth.dash}$, czyli około jednej dziesiątej. Wariancja
  po wygładzeniu to $\frac{1-\beta}{1+\beta}$ surowej, więc o obrazie decyduje
  ustawienie tak samo jak trening \dash{} Program~\ref{prog:P21} ze swoim
  własnym $\beta = \val{p33.beta}$ dostaje z tej samej formuły
  $\val{p33.smooth.pct}$ procent.}
\item Siedem punktów na $\val{p33.fit.decades}$ dekadach leży na prostej na
  osiach log-log, z dokładnością lepszą niż $\val{p33.fit.resid}$ procent. Co
  to wiąże w przewidywaniu dopasowania trzy dekady dalej?
  \teachesat{44--45}
  \answerto{Zupełnie nic. Reszta wiąże wewnątrz dopasowanego zakresu i nie
  mówi nic poza nim: trzy dekady dalej to samo dopasowanie myli się o
  $\val{p33.out.err}$ procent, a prawda leży poza jego własnym pasmem dwóch
  błędów standardowych.}
\end{quiz}

% ============================================================
\section{Jeden krok i sześć programów}
\label{sec:P33-one-step}
\index{trening!jeden krok}

\begin{fr}
Wszystko w tym programie zostało wyprowadzone gdzie indziej w tej książce. Taki
jest kontrakt Części IX i tutaj najłatwiej go sprawdzić, bo krok treningu jest
dość krótki, żeby go wypisać.

Zanim pójdziesz dalej: krok bierze partię przykładów i zwraca zmieniony zestaw
wag. Wymień to, co dzieje się pomiędzy \dash{} jest tego sześć i każde z tego
już spotkałeś.
\dotline
\nextframe
\end{fr}

\begin{fr}
\begin{ansblock}
Policzyć stratę. Zróżniczkować ją. Zamienić gradient w krok. Utrzymać
arytmetykę w zakresie reprezentowalnym. Zdecydować, jak daleko iść. I zrobić
to wszystko na \emph{próbce} danych, a nie na danych.
\end{ansblock}

Sześć elementów i żaden z nich nie jest tutaj nowy.
\end{fr}

\begin{fr}
Weź je po kolei i nazwij właściciela każdego.
\dotline
\nextframe
\end{fr}

\begin{fr}
\begin{ansblock}
\begin{itemize}
\item \textbf{Strata.} Program~\ref{prog:P26}: to ujemna logarytmiczna
  wiarygodność zaobserwowanych tokenów, więc trening modelu językowego jest
  metodą największej wiarygodności i nigdy nie był niczym innym. Sama liczba
  należy do Programu~\ref{prog:F02} i wynosi $\val{p33.loss.nats}$ natów.
\item \textbf{Jej gradient.} Program~\ref{prog:P18}: przez $\softmax$ i
  entropię krzyżową razem jest to $\vect{p} - \vect{y}$, a droga
  dwuetapowa, która liczy je osobno, jest inną funkcją, a nie wolniejszą.
\item \textbf{Krok.} Program~\ref{prog:P20} daje wzór, a w środku siedzi
  średnia wykładnicza Programu~\ref{prog:F04} wraz z jej korektą obciążenia.
\item \textbf{Arytmetyka.} Program~\ref{prog:P01}: odstęp między liczbami
  reprezentowalnymi rośnie z wielkością, więc wkład mniejszy niż połowa
  odstępu przy sumie bieżącej nie rusza niczego.
\item \textbf{Jak daleko.} Program~\ref{prog:P17}:
  $\eta < 2/\lambda_{\max}$, co jest twierdzeniem o ciągu geometrycznym i nie
  ma w nim nic o sieciach.
\item \textbf{Próbka.} Program~\ref{prog:P21}: gradient z jednej partii jest
  estymatorem, a prawie każda skarga na trening jest twierdzeniem o jego
  wariancji.
\end{itemize}
\end{ansblock}
\end{fr}

\mermaidfig{p33-six-pieces}{%
  Jeden krok jednego treningu, a obok każdego elementu program, który jest
  jego właścicielem. Część IX składa; nie wyprowadza.}{%
  krok i właściciel każdego elementu}

\begin{fr}
Dwa z nich niosą liczby, których ten program zaraz użyje, więc trafiają na
stronę teraz.

Granica Programu~\ref{prog:P17} na misie dzielonej z Programem~\ref{prog:P10}
to $\eta < \val{p33.eta.bound}$, a rozmiar kroku docierający najszybciej to
$\val{p33.eta.best}$. O jaką część różnią się najlepszy i zabójczy rozmiar
kroku?
\dotline
\nextframe
\end{fr}

\begin{fr}
\ans{$\val{p33.eta.margin}$ procent}
Tyle wynosi całe pasmo, w którym żyje harmonogram. Poniżej trening jest
wolniejszy, niż musi być; powyżej trening nie staje się gorszy, tylko odchodzi.

\begin{aibox}
To także powód, dla którego istnieje rozgrzewka współczynnika uczenia.
Wcześnie w treningu krzywizna nie jest tą krzywizną, pod którą strojono
harmonogram, więc to samo $\eta$ może wylądować po złej stronie granicy, która
się przesunęła. Program~\ref{prog:P20} zmierzył własny współczynnik rozgrzewki;
ta ramka dodaje, że chroniony margines ma szerokość kilku procent.
\end{aibox}
\end{fr}

\begin{fr}
Druga liczba to szum. Program~\ref{prog:P21} zmierzył rozrzut gradientu z
partii wprost; tutaj wielkością jest \emph{strata}, która jest średnią po
partii, więc tempo Programu~\ref{prog:P25} stosuje się bez zmian.

Weź stratę na przykład o rozrzucie $\val{p33.sigma.token}$ \dash{} to
najszerszy zmierzony wiersz Programu~\ref{prog:P19} \dash{} i partię
$\val{p33.batch}$. Jaki jest rozrzut liczby, którą wypisuje panel?
\dotline
\nextframe
\end{fr}

\begin{fr}
\ans{$\val{p33.sigma.step}$}
$\sigma/\sqrt{B}$, czyli sekcja 1 Programu~\ref{prog:P25} i nic więcej.

Zapamiętaj to. Każde twierdzenie w środku tego programu to ta liczba
przeciwko rozmiarowi rzeczy, którą próbujesz zobaczyć.
\end{fr}

% ============================================================
\section{Które z sześciu zawodzą po cichu}
\label{sec:P33-silent}
\index{trening!ciche awarie}

\begin{fr}
Teraz pytanie, które udostępnia dopiero złożenie i którego żaden z sześciu
właścicieli nie mógł zadać sam.

Zepsuj każdy z sześciu. Ile z nich zgłasza błąd, ostrzeżenie albo wyjątek?
\dotline
\nextframe
\end{fr}

\begin{fr}
\ans{$\val{p33.loud}$ z $\val{p33.pieces}$}
\begin{ansblock}
Uczciwa dwuetapowa droga gradientu tworzy $-1/p_c$, a gdy logit spadnie o
mniej więcej $\val{p33.grad.cliff}$ poniżej największego, $p_c$ znika do
dokładnego zera i droga dzieli przez zero. To jest głośne i
Program~\ref{prog:P18} to zmierzył.

Rozmiar kroku powyżej $\val{p33.eta.bound}$ mnoży najbardziej stromy kierunek
przez więcej niż jeden co krok. To też jest głośne, w końcu, a
Program~\ref{prog:P17} mówi dokładnie, gdzie leży granica.
\end{ansblock}

Nazwij pozostałe cztery i powiedz, ile kosztuje każdy z nich, gdy jest zły.
\dotline
\nextframe
\end{fr}

\begin{fr}
\begin{ansblock}
Strata, średnia w optymalizatorze, arytmetyka i próbka. Każde zmierzone przez
program, który jest jego właścicielem.
\end{ansblock}

\begin{center}
\begin{tabularx}{\linewidth}{@{}lX@{}}
\toprule
\textbf{Element} & \textbf{Ile kosztuje pomyłka i co ona mówi} \\
\midrule
strata & Uśrednij zakłopotania poszczególnych sekwencji zamiast
  eksponować średnią stratę, a raportowana liczba jest
  $\val{p33.fail.ppl}$ razy większa od prawdziwej przy najszerszym rozrzucie
  Programu~\ref{prog:P19}. \\
średnia & Przenieś rozmiar kroku między dwiema konwencjami pędu, a
  spacer Programu~\ref{prog:P20} liczący $\val{p33.conv.same}$ kroków trwa
  $\val{p33.conv.fold}$ razy dłużej albo rozbiega się. \\
arytmetyka & W \code{bf16} wkład poniżej $\val{p33.fail.swamp}$ procent sumy
  bieżącej nie rusza niczego, więc gradient można dodać i nie mieć żadnego
  efektu. \\
próbka & Akumuluj mikropartie, z których każda podzieliła już przez własną
  liczbę tokenów, a gradient dotyczy innego celu, o czynnik
  $\val{p33.fail.acc}$. \\
\bottomrule
\end{tabularx}
\end{center}
\end{fr}

\begin{fr}
\begin{trapbox}
\textbf{\enquote{Gdyby cokolwiek z tego się działo, framework by mi
powiedział.}} Powie ci o $\val{p33.loud}$ z $\val{p33.pieces}$, a tylko o
jednym z tych dwóch od razu.

Powód nie jest taki, że narzędzia są niedbałe. Każda z czterech cichych awarii
jest \emph{poprawnym obliczeniem złej wielkości}: średnia średnich jest
średnią, zaokrąglona suma jest sumą, przeskalowany krok jest krokiem. Nie ma w
nich nigdzie niedozwolonej operacji, którą cokolwiek mogłoby złapać.
\end{trapbox}
\end{fr}

\begin{fr}
Więc te cztery, które zawodzą po cichu, to te cztery, które trzeba wykryć z
zewnątrz, a narzędzie do tego jest dokładnie jedno.

Jakie?
\dotline
\nextframe
\end{fr}

\begin{fr}
\ans{krzywa straty}
Która jest jedną liczbą na krok, a \emph{każda z tych liczb ma rozrzut
$\val{p33.sigma.step}$ wokół tego, co mierzy}.

To jest cała reszta tego programu. Cztery kształty, które ludzie odczytują z
krzywej straty, to plateau, skok, podejrzanie gładki spadek i rozbieżność, a
pytanie zadawane każdemu z nich jest tym samym pytaniem:
\textbf{czy to jest dowód, czy to jest to, co robi szum?}
\end{fr}

% ============================================================
\section{Plateau}
\label{sec:P33-plateau}
\index{trening!plateau}

\mermaidfig{p33-flat-or-stuck}{%
  Ta sama obserwacja i dwie rzeczy, którymi może być. Reszta tej sekcji to
  jedna arytmetyka, która je rozróżnia \dash{} albo mówi, że w tej skali nie
  rozróżnia ich nic.}{%
  płaski odcinek i czym może być}

\begin{fr}
Obserwujesz od $\val{p33.watch}$ kroków, a strata na panelu teraz nie jest
niższa od odczytu, który pamiętasz sprzed $\val{p33.watch}$ kroków.

Załóż, że trening jest całkiem zdrowy i strata naprawdę spada o
$\val{p33.improve}$ na krok. Jak prawdopodobne jest, że zobaczyłbyś dokładnie
to?
\dotline
\nextframe
\end{fr}

\begin{fr}
\begin{ansblock}
Prawdziwy spadek przez $\val{p33.watch}$ kroków to $\val{p33.watch}$ razy
$\val{p33.improve}$. Każdy z dwóch odczytów niesie rozrzut
$\val{p33.sigma.step}$, więc ich różnica niesie
$\val{p33.sigma.step}\sqrt{2}$, a szansa braku widocznej poprawy to jeden
minus dystrybuanta rozkładu normalnego w ich ilorazie.

Wynosi $\val{p33.p.flat}$ procent.
\end{ansblock}

Prawie połowa przypadków. Zdrowy trening wygląda na zablokowany przez
$\val{p33.watch}$ kroków mniej więcej tak często, jak moneta wypada orłem.

\begin{note}
\textbf{To dwa odczyty i to jest pytanie, które warto zadać. Surowsze
odczytanie tego samego zdania daje zdarzenie o dwa rzędy wielkości rzadsze.}
\enquote{Strata nie zeszła poniżej wartości początkowej} może znaczyć, że
żaden odczyt na całym odcinku nigdy nie spadł pod pierwszy, a to
$\val{p33.watch}$ szans na porażkę, nie jedna. Policzone przy tym samym
modelu szumu \dash{} z warunkowaniem po wspólnym pierwszym odczycie, bo
porównania nie są od niego niezależne \dash{} daje to
$\val{p33.p.flat.min}$ procent przeciw $\val{p33.p.flat}$, czyli czynnik
$\val{p33.p.flat.ratio}$.

Obie liczby są prawdziwe o swoim własnym zdarzeniu i tylko jedna z nich mówi o
człowieku. Nikt nie przegląda $\val{p33.watch}$ zapisanych wartości i nie
bierze z nich minimum; patrzy się na liczbę teraz i pamięta, gdzie była z
grubsza wcześniej, czyli robi to porównanie powyżej \dash{} i z niego wyprowadza
się próg z tej sekcji. Surowsza liczba jest też bardziej krucha: potrzebuje
niezależności szumu krok po kroku, a prawdziwa zapisana strata jest
skorelowana między krokami, bo strumień partii przetasowano raz.
\end{note}

Jedna liczba w tej arytmetyce została zmierzona, a jedna wybrana. Która jest
która?
\dotline
\nextframe
\end{fr}

\begin{fr}
\begin{ansblock}
Rozrzut zmierzono; spadek wybrano.
\end{ansblock}

Warto powiedzieć dokładnie, jak daleko sięga każde z nich. Spadek o
$\val{p33.drop}$ przez $\val{p33.drop.steps}$ kroków to wybór zrobiony tutaj,
żeby mieć przez co dzielić. Rozrzut wyborem nie jest \dash{} to zmierzony przez
Program~\ref{prog:P19} rozrzut na token, któremu Program~\ref{prog:P25}
dokłada swoje $\sqrt{B}$.

To, co mówi arytmetyka, jest \emph{stosunkiem}, i ten stosunek jest sednem:
postęp na krok jest znacznie mniejszy niż szum na krok, a wszystko poniżej
wynika z tego samego.
\end{fr}

\begin{fr}
Odwróć więc pytanie, bo to jego jedyna użyteczna postać.

Ile kroków musi minąć bez poprawy, zanim płaskość stanie się dowodem \dash{}
zanim zdrowy trening pokazywałby ją rzadziej niż w $\val{p33.alpha}$
przypadków?
\dotline
\nextframe
\end{fr}

\begin{fr}
\ans{około $\val{p33.plateau.k}$ kroków}
\begin{ansblock}
Zestaw $k \times \val{p33.improve}$ z
$\val{p33.z.alpha} \times \val{p33.sigma.step}\sqrt{2}$ i rozwiąż względem
$k$.
\end{ansblock}

Trening w tym scenariuszu ma $\val{p33.drop.steps}$ kroków od początku do
końca. \textbf{Płaski odcinek musiałby być mniej więcej trzy razy dłuższy niż
cały trening, żeby cokolwiek znaczyć.}
\end{fr}

\transcript{p33-how-flat}

\begin{fr}
\begin{trapbox}
\textbf{\enquote{Strata jest płaska od godziny; coś jest nie tak.}} Przy tym
stosunku postępu do szumu godzina płaskości jest oczekiwanym zachowaniem
działającego treningu i żadne wpatrywanie się w krzywą tych dwóch nie
rozdzieli.

Błędem nie jest niecierpliwość. Błędem jest zadawanie pytania narzędziu, które
nie może na nie odpowiedzieć \dash{} a lekarstwem nie jest dłuższe patrzenie,
tylko inny pomiar: rzadsze, większe ewaluacje, gdzie Program~\ref{prog:P25}
i jego $1/\sqrt{n}$ pracują dla ciebie, a nie przeciwko tobie.
\end{trapbox}

Ta liczba jest dość duża, żeby warto ją było zaatakować. Podaj dwa powody, dla
których może być za duża.
\dotline
\nextframe
\end{fr}

\begin{fr}
\begin{ansblock}
Wczesny trening, gdzie postęp na krok jest większy; oraz szum dwóch odczytów,
który nie jest całkiem niezależny.
\end{ansblock}

Oba warto rozwinąć. Pierwszy: późny trening jest z konstrukcji trudnym
przypadkiem. Wcześnie strata spada szybko, postęp na krok jest duży przy tym
samym szumie, a plateau kilkuset kroków naprawdę jest informacją. Krzywa
przestaje móc powiedzieć cokolwiek dokładnie wtedy, gdy trening stał się
drogi.

Drugi: arytmetyka zakłada, że szum dwóch odczytów jest niezależny. Nie całkiem
\dash{} kolejne partie zachodzą na siebie w tym, co mierzą. To czyni rozrzut
różnicy nieco mniejszym niż $\val{p33.sigma.step}\sqrt{2}$, a odpowiedź nieco
mniejszą niż $\val{p33.plateau.k}$ \dash{} mniejszą, a nie innego rodzaju.
\end{fr}

% ============================================================
\section{Skok}
\label{sec:P33-spike}
\index{trening!skok}

\begin{fr}
Drugi kształt. Strata jednego kroku leży cztery rozrzuty powyżej trendu, a
wszystko po obu stronach jest zwyczajne.

Zanim uznasz to za zdarzenie: trening ma $\val{p33.steps}$ kroków. Ilu takich
kroków się spodziewałeś?
\dotline
\nextframe
\end{fr}

\begin{fr}
\ans{około $\val{p33.spike.expected}$}
Odchylenie o cztery rozrzuty ma prawdopodobieństwo mniej więcej trzy na sto
tysięcy, a kroków jest $\val{p33.steps}$. Pomnóż.

\textbf{Trzy z nich to oczekiwana liczba w czystym treningu.} Znalezienie
jednego nie jest znalezieniem niczego.
\end{fr}

\begin{fr}
Spotkałeś dokładnie tę arytmetykę wcześniej, jedną część temu i o modelach
zamiast o krokach.

Gdzie?
\dotline
\nextframe
\end{fr}

\begin{fr}
\ans{Program~\ref{prog:P27}}
Ten program zmierzył, co się dzieje, gdy czterdzieści modeli porównuje się na
jednej tablicy wyników i zwycięzcę wybiera się po fakcie: największy z wielu
losowań jest duży z konstrukcji, a próg trzeba skorygować o to, ile ich
wzięto.

Trening jest tym samym doświadczeniem z $\val{p33.steps}$ losowaniami zamiast
czterdziestu. Korekta o to daje rozmiar, który skok musi osiągnąć, zanim
stanie się zaskakujący:
\[ z = \val{p33.spike.z} \]
a skrypt, który to liczy, odtwarza zatwierdzoną liczbę
Programu~\ref{prog:P27}, gdy w miejsce $\val{p33.steps}$ wstawić czterdzieści,
więc te dwa są jednym rachunkiem, a nie dwoma podobnymi do siebie.

Dwa skoki tego samego rozmiaru. Jeden zostaje wchłonięty i krzywa idzie dalej
tam, gdzie była; drugi zostawia poziom wyżej. Który warto zbadać?
\dotline
\nextframe
\end{fr}

\begin{fr}
\ans{drugi}
\begin{aibox}
Skok, \emph{po którym następuje powrót}, jest jako sygnał wart prawie nic, a
skok bez powrotu jest wart bardzo dużo. Powrót jest informacją: odchylenie,
które trening wchłonął, było losowaniem z ogona, a odchylenie zmieniające
potem poziom było zdarzeniem.

Ta sama logika wycenia alarmy. Próg ustawiony na czterech rozrzutach w
treningu liczącym $\val{p33.steps}$ kroków budzi kogoś mniej więcej trzy razy
na trening bez powodu, a zespół, który nauczył się go ignorować, nauczył się
ignorować ten czwarty.
\end{aibox}
\end{fr}

\begin{fr}
Za ogonem stoi też mechanizm, który dał ci już Program~\ref{prog:F06}. Skok
\emph{straty} i skok \emph{gradientu} to różne obserwacje, a przycinanie
działa na drugim.

Program~\ref{prog:P21} zmierzył, ile przycina próg na czterokrotności typowego
rozmiaru gradientu: jeden krok na $\val{p33.steps}$ w jego własnej próbie, co
jest tą samą arytmetyką jeszcze raz. To, co robi przycinanie na połowie tego
progu, nie jest siatką bezpieczeństwa \dash{} to inny algorytm, bo zmienia
długość większości kroków, a nie kilku.
\end{fr}

% ============================================================
\section{Podejrzanie gładki spadek}
\label{sec:P33-smooth}
\index{trening!wygładzanie}

\begin{fr}
Trzeci kształt i to ten, którego nikt nie podejrzewa, bo gładka krzywa wygląda
na dobrą wiadomość.

Sekcja 5 Programu~\ref{prog:P21} ma połowę tego, co robi wygładzanie:
przesuwa każde zdarzenie później dokładnie o wybrany czas połowicznego zaniku,
więc \emph{kiedy} coś się stało nie jest pytaniem, na które wygładzona krzywa
umie odpowiedzieć. Jej pułapka mówi: czytaj surową krzywą dla \emph{kiedy}, a
wygładzoną dla \emph{czy}.

To jest ta druga połowa. Co wygładzanie robi z \emph{czy}?
\dotline
\nextframe
\end{fr}

\begin{fr}
\begin{ansblock}
Zmniejsza rozrzut, o czynnik wybrany mimochodem.

Średnia wykładnicza niezależnego szumu ma wariancję
$\frac{1-\beta}{1+\beta}$ razy wariancja surowa \dash{} to suma szeregu
geometrycznego, czyli własny obiekt Programu~\ref{prog:F04}. Widoczny rozrzut
to więc
\[ \sqrt{\frac{1-\beta}{1+\beta}} \]
rozrzutu, który trening ma.
\end{ansblock}
\end{fr}

\begin{fr}
Program~\ref{prog:P21} ma własne $\beta = \val{p33.beta}$; wstaw je, a
obraz pokazuje $\val{p33.smooth.pct}$ procent prawdziwego rozrzutu.

Teraz domyślna wartość panelu. Co pokazuje $\beta = \val{p33.beta.dash}$?
\dotline
\nextframe
\end{fr}

\begin{fr}
\ans{$\val{p33.smooth.dash}$ tego}
Mniej więcej jedną dziesiątą. \textbf{Krzywa wygładzona przy
$\val{p33.beta.dash}$ zaniża rozrzut treningu krok po kroku o czynnik
$\val{p33.smooth.dash.times}$}, a ten czynnik jest własnością ustawienia, a
nie treningu.

Zauważ, od czego on \emph{nie} zależy. Nie od modelu, nie od rozmiaru partii,
nie od tego, jak długo trening trwał, nie od wielkości zmniejszanego rozrzutu
\dash{} tylko od $\beta$. Dwie krzywe wygładzone przy tym samym $\beta$
pozostają więc porównywalne, jakkolwiek różne byłyby oba treningi, a dwie
wygładzone przy różnym $\beta$ nie są porównywalne, jakkolwiek podobne byłyby
treningi.

Zatem: dwa treningi są narysowane obok siebie i jeden jest widocznie
spokojniejszy. Co musisz wiedzieć, zanim powiesz cokolwiek?
\dotline
\nextframe
\end{fr}

\begin{fr}
\ans{jak każdy z nich został wygładzony}
\begin{trapbox}
\textbf{\enquote{Ten trening jest dużo stabilniejszy niż poprzedni.}} Dwa
treningi są porównywalne pod względem gładkości tylko wtedy, gdy wygładzono je
identycznie, a stała wygładzania to jedyna liczba na panelu treningowym,
której nikt nie zapisuje w raporcie.

Gorzej, zwykle nie jest ona w ogóle wybierana: jest tym, na czym zostawiono
suwak wykresu, a suwak nie jest pomiarem.
\end{trapbox}

Czy wygładzanie jest więc błędem?
\dotline
\nextframe
\end{fr}

\begin{fr}
\ans{nie}
Nic z powyższego nie jest argumentem przeciwko wygładzaniu i warto powiedzieć,
do czego ono służy.

Oko źle uśrednia zaszumiony ciąg, a dobrze śledzi linię, więc wygładzona
krzywa jest właściwym narzędziem do pytania \emph{czy poziom się przesuwa}.
Nie jest narzędziem do \emph{jak zaszumiony jest ten trening} ani do
\emph{kiedy to się stało}, a na oba odpowiada pewnie i źle.

Odczytaj ustawienie z wykresu, zanim odczytasz wykres.
\end{fr}

\mermaidfig{p33-smooth-hides}{%
  Dwie rzeczy, które robi wygładzanie, a tylko jedna z nich jest gdziekolwiek
  zapisana. Opóźnienie to pomiar Programu P21; zmniejszenie rozrzutu to ten
  pomiar, i oba ustawia ta sama liczba.}{%
  wygładzanie, opóźnienie i rozrzut}

% ============================================================
\section{Jedyny kształt, który jest diagnozą}
\label{sec:P33-diverge}
\index{trening!rozbieżność}

\begin{fr}
Trzy kształty do tej pory i żaden nie przetrwał pytania. Czwarty przetrwa.

Strata rośnie i rośnie dalej. Który program tej książki mówi już dokładnie,
kiedy to się dzieje, i dlaczego nie potrzebuje żadnej statystyki?
\dotline
\nextframe
\end{fr}

\begin{fr}
\ans{Program~\ref{prog:P17}}
\begin{ansblock}
W modelu kwadratowym każdy kierunek własny jest mnożony przez
$1 - \eta\lambda$ co krok, więc spacer jest ograniczony dokładnie wtedy, gdy
$\lvert 1 - \eta\lambda \rvert < 1$ we wszystkich kierunkach naraz. Powyżej
$\eta = \val{p33.eta.bound}$ najbardziej stromy kierunek ma czynnik większy
niż jeden, a czynnik większy niż jeden stosowany co krok nie jest szumem.
\end{ansblock}

Program~\ref{prog:P17} ma własny pokazowy rozmiar kroku
$\val{p33.eta.over}$; przy nim czynnik wynosi $\val{p33.diverge.factor}$, a
odległość od minimum rośnie o tyle co krok na zawsze.

Cztery kształty zatem. Jak się dzielą?
\dotline
\nextframe
\end{fr}

\begin{fr}
\ans{dwa na dwa, i nie tak, jak sugeruje panel}

\begin{center}
\begin{tabularx}{\linewidth}{@{}lX@{}}
\toprule
\textbf{Kształt} & \textbf{Czym jest} \\
\midrule
plateau & Prawie nigdy dowodem w skali, w jakiej obserwuje się trening. \\
skok & Oczekiwany, w liczbie dyktowanej przez długość treningu. \\
gładkość & Twierdzenie o wygładzaniu tak samo jak o treningu. \\
rozbieżność & Dowód, z progiem będącym własnością powierzchni. \\
\bottomrule
\end{tabularx}
\end{center}

Te trzy, które są czytane bez przerwy, mówią najmniej, a ten, który mówi
najwięcej, jest tym, którego nikt nie potrzebuje wykresu, żeby zauważyć.

Jeśli krzywa jest kiepskim narzędziem, jakie jest lekarstwo?
\dotline
\nextframe
\end{fr}

\begin{fr}
\ans{lepszy pomiar, a nie lepszy wykres}
\begin{aibox}
Krzywa jest kiepskim narzędziem, bo jest jedną zaszumioną liczbą na krok, więc
lekarstwem jest lepszy pomiar, a nie lepszy wykres, a
Program~\ref{prog:P27} daje obie jego połowy \dash{} ewaluuj na zbiorze
odłożonym, raportuj przedział zamiast punktu i paruj porównanie, żeby szum się
skrócił.

Ewaluacja co tysiąc kroków na ustalonym zbiorze odpowiada na pytania, na które
krzywa straty nie odpowie w żadnej rozdzielczości.
\end{aibox}
\end{fr}

% ============================================================
\section{Druga krzywa: prawo skalowania}
\label{sec:P33-scaling}
\index{prawo skalowania}

\begin{fr}
Drugi wykres w każdym opisie treningu wcale nie jest krzywą straty. To garść
skończonych treningów na osiach log-log z prostą przez nie i strzałką
wskazującą w prawo.

Program~\ref{prog:F03} dał ci maszynerię: prawo potęgowe $y = ax^{b}$ jest
prostą w $\ln x$ i $\ln y$, a nachylenie prostej jest wykładnikiem. Czego
celowo nie zrobił, to nie dopasował żadnego. Ta sekcja dopasowuje, a potem
pyta, ile takie dopasowanie jest warte.
\end{fr}

\begin{fr}
Oto doświadczenie, a wszystko w nim jest znane, bo ta książka to wybrała.
Prawdą jest
\[ L(N) = \val{p33.fit.linf} + a N^{-\val{p33.fit.alpha}} \]
czyli kształt używany przez samą literaturę: nieredukowalna strata, której
żadna skala nie usuwa, plus część, która się skaluje.

Bierze się z tego siedem punktów, równo rozłożonych na
$\val{p33.fit.decades}$ dekadach od $\val{p33.fit.lo}$ do
$\val{p33.fit.hi}$, a strata biegnie od $\val{p33.fit.loss.lo}$ w dół do
$\val{p33.fit.loss.hi}$. Do nich dopasowuje się \emph{czyste} prawo potęgowe
metodą najmniejszych kwadratów w $\ln N$ i $\ln L$, czyli to, co dopasowują
wszyscy.
\end{fr}

\begin{fr}
Pierwsze pytanie i to na nie używa się tego wykresu.

Prawdziwy wykładnik przy części redukowalnej to $\val{p33.fit.alpha}$. Jaki
wykładnik raportuje dopasowana prosta?
\dotline
\nextframe
\end{fr}

\begin{fr}
\ans{$\val{p33.fit.b}$}
Nie blisko i nie z winy dopasowania. Prosta jest najlepszą prostą przez tych
siedem punktów, a tych siedem punktów nie leży na prostej \dash{} leżą na
krzywej wyginanej przez podłogę. Dopasowany wykładnik jest własnością
\emph{zakresu} i wyszedłby inaczej na innym.

\textbf{Liczba odczytana z wykresu linijką nie jest więc wykładnikiem w tej
zależności.} To średnie nachylenie na zmierzonych dekadach.
\end{fr}

\begin{fr}
Drugie pytanie. Skoro tak, jak bardzo dopasowana prosta mija te siedem
punktów, do których ją dopasowano?
\dotline
\nextframe
\end{fr}

\begin{fr}
\ans{najwyżej o $\val{p33.fit.resid}$ procent}
Co jest niewidoczne. Przy szerokości narysowanej linii na
$\val{p33.fit.decades}$ dekadach pół procent nie jest rozbieżnością, którą
czytelnik zobaczy, i żadne wpatrywanie się w wykres nie zasugeruje, że
zależność jest czymkolwiek innym niż prawem potęgowym.

\textbf{Dopasowanie znakomite wszędzie tam, gdzie je sprawdzono, zwróciło
właśnie wykładnik zły prawie trzykrotnie.}
\end{fr}

\begin{fr}
Trzecie pytanie i to jego dotyczy strzałka na wykresie.

Ekstrapoluj dopasowaną prostą $\val{p33.fit.decades}$ dekady za jej ostatni
punkt, do $N = 10^{\val{p33.out.dec}}$. Prawda tam wynosi
$\val{p33.out.truth}$. Co przewiduje dopasowanie?
\dotline
\nextframe
\end{fr}

\begin{fr}
\ans{$\val{p33.out.fitted}$}
Za mało o $\val{p33.out.err}$ procent \dash{} przy reszcie wewnątrz zakresu
wynoszącej $\val{p33.fit.resid}$ procent.

A błąd nie jest ustalonym rozmiarem, który dałoby się podać jako margines.
Rośnie z każdą dekadą i rośnie bez ograniczenia, bo dopasowane prawo potęgowe
spada do zera, a prawda nie.

Prawda ma podłogę, a dopasowana prosta nie. Co to każe dopasowaniu przewidywać
dostatecznie daleko?
\dotline
\nextframe
\end{fr}

\begin{fr}
\ans{stratę poniżej podłogi}
I jest to wniosek, który można rozwiązać zamiast się o niego spierać.

Dopasowana prosta osiąga nieredukowalną stratę $\val{p33.fit.linf}$ przy
$N = 10^{\val{p33.absurd.dec}}$, czyli $\val{p33.absurd.past}$ dekady za
ostatnim zmierzonym punktem. Dalej przewiduje stratę \emph{poniżej} podłogi
\dash{} odpowiedź nie tylko niedokładną, ale niemożliwą, wyprodukowaną przez
dopasowanie o reszcie $\val{p33.fit.resid}$ procent.
\end{fr}

\mermaidfig{p33-inside-outside}{%
  Jedna prosta, dopasowana na trzech dekadach, i dwa różne twierdzenia, które
  ludzie z nią wypowiadają. W zakresie mierzona; poza nim zgadywana.}{%
  dopasowanie w zakresie i poza nim}

\begin{fr}
Zostaje ostatnie pytanie i o nie prosi kontrakt: dopasowanie raportuje własną
niepewność. Czego ta niepewność dotyczy?
\dotline
\nextframe
\end{fr}

\begin{fr}
\begin{ansblock}
Reszt \dash{} i niczego więcej.

Błąd standardowy nachylenia najmniejszych kwadratów jest zbudowany z tego, jak
daleko punkty mijają prostą. Te punkty mijają ją prawie o nic, więc
raportowany błąd standardowy wynosi $\val{p33.fit.se}$, a dwa takie
przeniesione przez $\val{p33.fit.decades}$ dekady dają pasmo od
$\val{p33.fit.band.lo}$ do $\val{p33.fit.band.hi}$.
\end{ansblock}

Prawda tam wynosi $\val{p33.out.truth}$, czyli leży poza tym pasmem.
\textbf{Dopasowanie jest pewne i błędne, a pewne jest dokładnie dlatego, że
jest błędne}: dane leżą na gładkiej krzywej, więc reszty są maleńkie, więc
raportowana niepewność jest maleńka.
\end{fr}

\begin{fr}
\begin{trapbox}
\textbf{\enquote{Prawo skalowania mówi, że dojdziemy do tej straty przy tym
budżecie.}} Dopasowana niepewność dotyczy szumu w pomiarach. Nie mówi nic o
tym, czy postać funkcyjna jest właściwa, a ekstrapolacja zależy wyłącznie od
postaci funkcyjnej.

Uczciwa forma twierdzenia nazywa swój zakres: \emph{w zmierzonym przedziale
strata spadała jak prawo potęgowe z tym wykładnikiem.} Wszystko na prawo od
ostatniego punktu jest założeniem o kształcie, a wykres nie zawiera o nim
żadnego dowodu \dash{} i właśnie dlatego rysuje się je z przedłużoną linią i
zaznaczonymi danymi.
\end{trapbox}
\end{fr}

\begin{fr}
\begin{aibox}
Istnieje praktyczna postać tego, która nic nie kosztuje. Dopasuj prawo dwa
razy \dash{} raz na wszystkich punktach, raz z odłożonym największym
treningiem \dash{} i porównaj, co to drugie przewiduje dla treningu, który już
masz.

To dyscyplina Programu~\ref{prog:P27} zastosowana do dopasowania, a nie do
tablicy wyników, i mierzy wielkość, która naprawdę się liczy: nie to, jak
dobrze prosta opisuje punkty, tylko jak dobrze prosta zbudowana z małych
treningów przewidziała ten duży.
\end{aibox}
\end{fr}

% ============================================================
\section{Czego ten program nie sprawdził}
\label{sec:P33-unchecked}
\index{pomiar!czego nie zmierzono}

\begin{fr}
Wszystko w sekcjach od 3 do 6 opiera się na jednym założeniu i warto
powiedzieć, na którym, bo reszta arytmetyki jest dokładna.

Nazwij je.
\dotline
\nextframe
\end{fr}

\begin{fr}
\begin{ansblock}
Że odchylenie straty krok po kroku ma rozkład normalny o rozrzucie, który się
nie zmienia.
\end{ansblock}

Pierwsza połowa to centralne twierdzenie graniczne Programu~\ref{prog:P25}
robiące swoje \dash{} średnia z partii wielu przykładów jest bliska normalnej
w środku, a arytmetyka plateau używa środka. Ale Program~\ref{prog:P25}
zmierzył też, gdzie to przestaje być prawdą: przy trzech rozrzutach rozkład
Gaussa już zawyża ogon, a arytmetyka skoku używa \emph{ogona}.

Wynik o plateau stoi więc na twardszym gruncie niż wynik o skoku i dlatego są
podane w tej kolejności.

Nie ma też obrony w kierunku błędu, bo ten kierunek nie jest jeden.
Program~\ref{prog:P25} zmierzył stosunek ogona rozkładu Gaussa do prawdy przy
jednym, trzech i pięciu rozrzutach: $\val{p25.tail.ratio1}$,
$\val{p25.tail.ratio3}$ i $\val{p25.tail.ratio5}$. Blisko środka
\emph{zaniża}, potem przecina prawdę i ucieka. Zdanie \enquote{to
przybliżenie jest ostrożne} nie jest więc czymś, co ktokolwiek może o nim
powiedzieć bez podania odległości.

A to, co by tę sprawę rozstrzygnęło, nic nie kosztuje i dlatego warto to
nazwać: zalogowane straty samego treningu \emph{są} próbką. Dopasuj trend,
weź reszty i zapytaj, jak często przekraczają trzy i cztery rozrzuty, w
zestawieniu z tym, co przewiduje rozkład normalny. Nikt tego nie robi, dane
już leżą na dysku, a odpowiedź zamieniłaby każdą liczbę z
\S\ref{sec:P33-spike} z modelu w pomiar.

To była pierwsza połowa założenia. Która połowa jest słabsza?
\dotline
\nextframe
\end{fr}

\begin{fr}
\ans{druga}
Szum prawdziwego treningu nie jest stacjonarny: rozrzut zmienia się wraz z
modelem, kolejność danych nie jest losowa w sposób, który zakłada arytmetyka,
a skok jest często własnością konkretnej partii, a nie losowaniem z ustalonego
rozkładu.

\begin{warning}
\textbf{Ta książka nie zmierzyła nic z tego i potrzebowałaby do tego
prawdziwego treningu.} To, co ustalają powyższe sekcje, jest \emph{podłogą}:
nawet przy założeniach łaskawych dla krzywej plateau jest nieinformatywne
przez dziesiątki tysięcy kroków, a długi trening zawiera duże odchylenia z
konstrukcji. Prawdziwy szum, cięższy w ogonach i dryfujący, pogarsza oba, a
nie poprawia.
\end{warning}
\end{fr}

\begin{fr}
I to jest uczciwe stanowisko Części IX, osiągnięte już trzy razy.
Program~\ref{prog:P32} podzielił własne zaległe twierdzenie tak samo: połowa o
złożeniu nie potrzebowała wytrenowanego modelu, a połowa o treningu
potrzebowała, więc pierwszą wyprowadzono, a drugą podano jako zaległą.

Zadaj ten podział wszystkiemu w tej książce, co wygląda, jakby potrzebowało
modelu. Zaskakująco dużo z tego okazuje się twierdzeniem o kształtach, o
rozkładzie normalnym albo o prostej najmniejszych kwadratów \dash{} a ta
część, która naprawdę potrzebuje treningu, jest warta nazwania zamiast
osłabiania całego twierdzenia.
\end{fr}

\begin{summarybox}
\sumitem{1--4}{Krok treningu to sześć elementów i ta książka ma każdy z nich:
  strata (Program~\ref{prog:P26}), jej gradient
  (Program~\ref{prog:P18}), aktualizacja (Program~\ref{prog:P20}),
  arytmetyka (Program~\ref{prog:P01}), rozmiar kroku
  (Program~\ref{prog:P17}) i próbka (Program~\ref{prog:P21}).}
\sumitem{5--6}{\result{Najlepszy rozmiar kroku na misie Programu~\ref{prog:P17} to
  $\val{p33.eta.best}$, a zabójczy $\val{p33.eta.bound}$: pasmo szerokie na
  $\val{p33.eta.margin}$ procent}.}
\sumitem{7--8}{\result{Średnia z partii straty o rozrzucie na przykład
  $\val{p33.sigma.token}$ po $\val{p33.batch}$ przykładach ma rozrzut
  $\val{p33.sigma.step}$}.}
\sumitem{9--11}{\result{Z sześciu elementów $\val{p33.loud}$ zgłaszają coś, a
  $\val{p33.silent}$ nie, bo każdy z tych czterech jest poprawnym
  obliczeniem złej wielkości}.}
\sumitem{12--14}{Te cztery ciche są wykrywalne tylko z zewnątrz, a narzędziem
  jest krzywa straty zbudowana z zaszumionych liczb.}
\sumitem{15--17}{\result{Zdrowy trening, którego strata spada o $\val{p33.improve}$ na
  krok, nie pokazuje poprawy między dwoma odczytami odległymi o
  $\val{p33.watch}$ kroków w $\val{p33.p.flat}$ procentach przypadków}
  \dash{} a surowsze odczytanie, przy którym żaden odczyt na odcinku nigdy nie
  spadł pod pierwszy, to $\val{p33.p.flat.min}$ procent, czyli
  $\val{p33.p.flat.ratio}$ razy rzadziej.}
\sumitem{18--21}{\result{Płaski odcinek musi trwać około $\val{p33.plateau.k}$ kroków,
  zanim zdrowy trening produkowałby go rzadziej niż w $\val{p33.alpha}$
  przypadków}.}
\sumitem{22--23}{\result{Trening liczący $\val{p33.steps}$ kroków zawiera około
  $\val{p33.spike.expected}$ odchyleń o cztery rozrzuty niezależnie od tego,
  co jeszcze jest o nim prawdą}.}
\sumitem{24--27}{\result{Po korekcie o długość treningu skok musi przekroczyć
  $\val{p33.spike.z}$ rozrzutu, żeby zaskakiwał} \dash{} arytmetyka
  Programu~\ref{prog:P27} z krokami w miejsce modeli.}
\sumitem{28--33}{\result{Średnia wykładnicza pokazuje
  $\sqrt{(1-\beta)/(1+\beta)}$ prawdziwego rozrzutu:
  $\val{p33.smooth.pct}$ procent przy $\beta = \val{p33.beta}$ i czynnik
  $\val{p33.smooth.dash.times}$ przy $\val{p33.beta.dash}$}.}
\sumitem{34--37}{\result{Rozbieżność to jedyny kształt niepotrzebujący statystyki:
  powyżej $\eta = \val{p33.eta.bound}$ czynnik większy niż jeden jest
  stosowany co krok}.}
\sumitem{40--43}{\result{Czyste prawo potęgowe dopasowane do prawdy postaci
  podłoga-plus-prawo-potęgowe mija własne siedem punktów o
  $\val{p33.fit.resid}$ procent i zwraca wykładnik $\val{p33.fit.b}$, gdzie
  prawda to $\val{p33.fit.alpha}$}.}
\sumitem{44--47}{To samo dopasowanie jest za małe o $\val{p33.out.err}$
  procent $\val{p33.fit.decades}$ dekady dalej i przewiduje stratę poniżej
  nieredukowalnej podłogi przy $10^{\val{p33.absurd.dec}}$.}
\sumitem{48--50}{\result{Jego własna raportowana niepewność jest zbudowana z reszt,
  więc jest mała dokładnie wtedy, gdy postać funkcyjna jest zła, a prawda leży
  poza jego pasmem}.}
\sumitem{51--54}{Nic z tego nie zostało sprawdzone na szumie prawdziwego
  treningu, który ma cięższe ogony i dryfuje; powyższe wyniki są podłogą, a
  nie oszacowaniem.}
\end{summarybox}

\begin{testexercises}
\item Strata treningu spada o $\num{2e-5}$ na krok, a obserwowana strata ma
  rozrzut $\num{0.2}$. Mniej więcej ile kroków bez poprawy produkowałby zdrowy
  trening pięć razy na sto?
  \answerto{Około $\num{23000}$. Próg to
  $\val{p33.z.alpha} \times \num{0.2}\sqrt{2}$ podzielone przez $\num{2e-5}$.
  Zauważ, że podwojenie postępu na krok zmniejszyło go o połowę, a
  zmniejszenie rozrzutu o połowę zrobiłoby to samo \dash{} odpowiedź zależy od
  stosunku i od niczego więcej.}
\item Zespół podnosi rozmiar partii z $\val{p33.batch}$ do $\num{1024}$ i
  raportuje, że krzywa straty stała się dużo czystsza. O jaki czynnik spadł
  rozrzut i czy trening się poprawił?
  \answerto{O $2$, bo rozrzut to $\sigma/\sqrt{B}$, a partia wzrosła
  czterokrotnie. Trening niekoniecznie poprawił się w ogóle: czystsza krzywa
  przy większej partii to dokładnie to, co przewiduje tempo
  Programu~\ref{prog:P25}, cokolwiek model robi.}
\item Trening liczący $\num{200000}$ kroków jest monitorowany alarmem przy
  trzech rozrzutach powyżej trendu. Ilu alarmów należy się spodziewać w
  czystym treningu?
  \answerto{Około $\num{270}$. Odchylenie o trzy rozrzuty ma
  prawdopodobieństwo mniej więcej $\num{0.00135}$, a $\num{200000}$ kroków je
  mnoży. Alarm odzywający się setki razy na trening nie jest alarmem.}
\item Rysowane są dwa treningi, jeden wygładzony przy $\beta = \num{0.9}$, a
  drugi przy $\beta = \num{0.99}$. Ich narysowane krzywe mają widoczne
  rozrzuty w stosunku $3$ do $1$. Jaki jest stosunek rozrzutów leżących pod
  spodem treningów?
  \answerto{Policz dwa czynniki zmniejszające i podziel je:
  \[ \sqrt{\num{0.1}/\num{1.9}} = \num{0.229},
     \qquad \sqrt{\num{0.01}/\num{1.99}} = \num{0.0709} \]
  czyli stosunek
  $\num{3.24}$. Narysowany stosunek $3$ jest więc bardzo bliski temu, co
  dałyby identyczne treningi, a dowodów na jakąkolwiek różnicę między nimi
  jest w przybliżeniu tyle co nic.}
\item Prawo potęgowe dopasowane na dwóch dekadach ma resztę jednego procenta.
  Co to wiąże w jego przewidywaniu jedną dekadę dalej?
  \answerto{Zupełnie nic. Reszta wiąże dopasowanie wewnątrz zakresu, na
  którym je dopasowano. Poza nim błąd zależy od tego, czy postać funkcyjna
  jest właściwa, a reszta nie niesie o tym żadnej informacji \dash{} co
  pokazuje własne dopasowanie tego programu przy
  $\val{p33.fit.resid}$ procent.}
\end{testexercises}

\canyou

\begin{furtherproblems}
\item Pokaż, że jeśli obserwowana strata w każdym kroku jest prawdą plus
  niezależny szum o rozrzucie $\sigma$, to różnica między odczytami odległymi
  o $k$ kroków ma rozrzut $\sigma\sqrt{2}$ niezależnie od $k$.
  \answerto{Oba składniki szumu są niezależne, więc ich wariancje się dodają:
  $\sigma^{2} + \sigma^{2}$. To $k$ wchodzi w \emph{prawdziwy} spadek, a nie w
  szum, i właśnie dlatego próg plateau jest liniowy w $k$.}
\item Zespół akumuluje cztery mikropartie o $1000$, $10$, $500$ i $490$
  prawdziwych tokenach, z których każda podzieliła już przez własną liczbę.
  Ich średnie straty to $2$, $8$, $3$ i $3$. Czy zakumulowany gradient jest
  tym z dużej partii?
  \answerto{Nie. Średnia ważona tokenami daje $\val{p33.accum.pooled}$;
  zakumulowana uśrednia te cztery i daje $\val{p33.accum.naive}$, bo partia
  dziesięciotokenowa niesie jedną dwusetną tokenów i ćwierć wagi. Błąd jest
  cichy, jest w gradiencie, a nie w raporcie, i jest pułapką
  Programu~\ref{prog:P21} z jedną mikropartią więcej.}
\item Średnia wykładnicza jest zastosowana dwa razy, przy tym samym $\beta$.
  O jaki czynnik zmniejszony jest rozrzut podwójnie wygładzonej krzywej?
  \answerto{Mniej niż o kwadrat pojedynczego czynnika, bo ciąg raz wygładzony
  nie jest już niezależny krok po kroku. Wart zabrania jest kierunek: każde
  przejście czyni obraz spokojniejszym, a trening ani trochę innym, i dwa
  przejścia nie są dwa razy uczciwsze.}
\item Dopasowanie z \S\ref{sec:P33-scaling} osiąga nieredukowalną stratę przy
  $10^{\val{p33.absurd.dec}}$. Rozwiąż to równanie względem tego rozmiaru z
  dopasowanej prostej.
  \answerto{Przyrównaj dopasowane $\exp(c) N^{-b}$ do
  $\val{p33.fit.linf}$ i weź logarytmy:
  $\ln N = (c - \ln \val{p33.fit.linf})/b$. Arytmetyka jest przez cały czas
  Programu~\ref{prog:F03} i to jest sedno \dash{} absurd widać przy pomocy
  niczego więcej niż dwóch praw logarytmów.}
\item Kolega proponuje wykrywanie plateau przez wcześniejsze mocne wygładzenie
  krzywej, na tej podstawie, że wygładzona krzywa jest mniej zaszumiona. Co
  jest z tym nie tak?
  \answerto{Wygładzona krzywa jest mniej zaszumiona, a oszacowanie, z którego
  jest zbudowana, nie. Wygładzanie zmniejsza \emph{widoczny} rozrzut o znany
  czynnik, nie dodając żadnej informacji, więc czyni plateau łatwiejszym do
  zobaczenia i ani trochę łatwiejszym do uwierzenia \dash{} a moment, w którym
  je zobaczysz, opóźnia o czas połowicznego zaniku. Zmienić trzeba narzędzie,
  a nie wykres.}
\end{furtherproblems}
